\section*{Chapitre \ref{ch:g12:complex} --- Nombres complexes}

\begin{solution}{exo:g12:complex:1}
$(3-2\iu)(1+4\iu) = 3 + 12\iu - 2\iu - 8\iu^2 = 11 + 10\iu$.

$\dfrac{2+\iu}{1-\iu} = \dfrac{(2+\iu)(1+\iu)}{(1-\iu)(1+\iu)}
= \dfrac{1 + 3\iu}{2} = \dfrac12 + \dfrac32\iu$.

$\iu^{2027}$~: comme $\iu^4 = 1$ et $2027 = 4 \times 506 + 3$,
$\iu^{2027} = \iu^3 = -\iu$.
\end{solution}

\begin{solution}{exo:g12:complex:2}
\emph{(a)} $\Delta = 16 - 52 = -36$~:
$z = \dfrac{4 \pm 6\iu}{2} = 2 \pm 3\iu$. Paire conjuguée~; produit
$= 4 + 9 = 13 = \frac ca$. \emph{(b)} $\Delta = 1 - 4 = -3$~:
$z = \dfrac{-1 \pm \iu\sqrt3}{2}$ (les racines cubiques de l'unité $j$ et
$\conj j$)~; produit $= \frac{1 + 3}{4} = 1$.
\end{solution}

\begin{solution}{exo:g12:complex:3}
$\abs{z_1} = \sqrt2$, \omterm{def:g12:complex:argument}{argument} $\frac{3\pi}{4}$ (deuxième quadrant)~:
$z_1 = \sqrt2\,\eu^{3\iu\pi/4}$.
$\abs{z_2} = 2$, \omterm{def:g12:complex:argument}{argument} $-\frac\pi6$~:
$z_2 = 2\,\eu^{-\iu\pi/6}$. D'où
\[
\frac{z_1}{z_2} = \frac{\sqrt2}{2}\,\eu^{\iu\left(\frac{3\pi}{4} + \frac{\pi}{6}\right)}
= \frac{\sqrt2}{2}\,\eu^{11\iu\pi/12}.
\]
Algébriquement,
\[
\frac{z_1}{z_2} = \frac{-1+\iu}{\sqrt3 - \iu}
= \frac{(-1+\iu)(\sqrt3+\iu)}{4}
= \frac{-\sqrt3 - 1 + \iu(\sqrt3 - 1)}{4}.
\]
En identifiant les \omterm{def:g12:complex:def}{parties réelles}~:
$\frac{\sqrt2}{2}\cos\frac{11\pi}{12} = \frac{-\sqrt3-1}{4}$, donc
\[
\cos\frac{11\pi}{12} = -\frac{\sqrt6 + \sqrt2}{4}.
\]
\end{solution}

\begin{solution}{exo:g12:complex:4}
\emph{(a)} $\abs{z - 2} = \abs{z - (-\iu)}$ signifie que $M$ est
équidistant de $A(2, 0)$ et $B(0, -1)$~: la \emph{médiatrice} de $[AB]$.

\emph{(b)} Distance $3$ du point $(1,1)$~: le \emph{cercle} de centre
$1 + \iu$ et de rayon $3$.

\emph{(c)} Par la \cref{prop:g12:complex:geometry}(4), $\frac{z-1}{z+1}$
purement imaginaire signifie que les \omterm{def:g11:vect:vector}{vecteurs} de $A(1,0)$ vers $M$ et de
$B(-1,0)$ vers $M$ sont \omterm{def:g11:scal:orthogonal}{orthogonaux}~: $M$ voit le segment $[AB]$ sous un
angle droit. L'ensemble est le \emph{cercle de diamètre $[AB]$} (le
cercle unité) \emph{privé des points $A$ et $B$} eux-mêmes.
\end{solution}

\begin{solution}{exo:g12:complex:5}
Moivre~: $\cos3\theta + \iu\sin3\theta = (\cos\theta + \iu\sin\theta)^3$.
En développant par la formule du binôme avec $c = \cos\theta$,
$s = \sin\theta$~:
\[
(c + \iu s)^3 = c^3 + 3\iu c^2 s - 3 c s^2 - \iu s^3
= (c^3 - 3cs^2) + \iu(3c^2 s - s^3).
\]
En identifiant les parties et en utilisant $s^2 = 1 - c^2$,
$c^2 = 1 - s^2$~:
\[
\cos3\theta = 4\cos^3\theta - 3\cos\theta,
\qquad
\sin3\theta = 3\sin\theta - 4\sin^3\theta .
\]
\end{solution}

\begin{solution}{exo:g12:complex:6}
\[
\cos^3\theta = \left(\frac{\eu^{\iu\theta} + \eu^{-\iu\theta}}{2}\right)^{\!3}
= \frac{\eu^{3\iu\theta} + 3\eu^{\iu\theta} + 3\eu^{-\iu\theta} + \eu^{-3\iu\theta}}{8}
= \frac{\cos3\theta + 3\cos\theta}{4}.
\]
D'où
\[
\int_0^{\pi/2}\cos^3\theta\,\dd\theta
= \frac14\left[\frac{\sin3\theta}{3} + 3\sin\theta\right]_0^{\pi/2}
= \frac14\left(-\frac13 + 3\right) = \frac{2}{3}.
\]
(Cohérent avec $W_3 = \frac23$ dans l'\cref{exo:g12:integ:9}, car
$\int_0^{\pi/2}\cos^3 = \int_0^{\pi/2}\sin^3$ par la symétrie
$\theta \mapsto \frac\pi2 - \theta$.)
\end{solution}

\begin{solution}{exo:g12:complex:7}
$c = a + \eu^{\pm\iu\pi/3}(b - a)$ avec $b - a = 2 + \iu$ et
$\eu^{\pm\iu\pi/3} = \frac12 \pm \iu\frac{\sqrt3}{2}$~:
\[
\eu^{\iu\pi/3}(2+\iu) = \left(\tfrac12 + \iu\tfrac{\sqrt3}{2}\right)(2+\iu)
= \left(1 - \tfrac{\sqrt3}{2}\right) + \iu\left(\tfrac12 + \sqrt3\right),
\]
donc $c_1 = \left(2 - \frac{\sqrt3}{2}\right) +
\iu\left(\frac32 + \sqrt3\right)$, et de même avec l'angle
$-\frac\pi3$~:
$c_2 = \left(2 + \frac{\sqrt3}{2}\right) + \iu\left(\frac32 - \sqrt3\right)$.
\end{solution}

\begin{solution}{exo:g12:complex:8}
$-4 = 4\,\eu^{\iu\pi}$, donc les solutions de $z^4 = -4$ sont
\[
z_k = \sqrt2\,\eu^{\iu\left(\frac\pi4 + \frac{k\pi}{2}\right)},
\quad k = 0, 1, 2, 3,
\]
\emph{c.-à-d.} $1 + \iu$, $-1 + \iu$, $-1 - \iu$, $1 - \iu$~: les sommets
d'un carré de circonférence de rayon $\sqrt 2$. En groupant les \omterm{def:g11:quad:discriminant}{racines}
\omterm{def:g12:complex:conjugate}{conjuguées},
\[
z^4 + 4 = \bigl(z - (1{+}\iu)\bigr)\bigl(z - (1{-}\iu)\bigr)
\bigl(z + (1{-}\iu)\bigr)\bigl(z + (1{+}\iu)\bigr)
= (z^2 - 2z + 2)(z^2 + 2z + 2).
\]
\end{solution}

\begin{solution}{exo:g12:complex:9}
Avec $q = \eu^{\iu\theta} \neq 1$~:
\[
\sum_{k=0}^{n} \eu^{\iu k\theta} = \frac{\eu^{\iu(n+1)\theta} - 1}{\eu^{\iu\theta} - 1}
= \frac{\eu^{\iu(n+1)\theta/2}}{\eu^{\iu\theta/2}}
\cdot\frac{\eu^{\iu(n+1)\theta/2} - \eu^{-\iu(n+1)\theta/2}}{\eu^{\iu\theta/2} - \eu^{-\iu\theta/2}}
= \eu^{\iu n\theta/2}\,
\frac{\sin\frac{(n+1)\theta}{2}}{\sin\frac{\theta}{2}},
\]
en utilisant $\eu^{\iu\alpha} - \eu^{-\iu\alpha} = 2\iu\sin\alpha$ (le
\emph{factorisation en demi-angle}). En prenant les \omterm{def:g12:complex:def}{parties réelles}~:
\[
S = \frac{\sin\frac{(n+1)\theta}{2}}{\sin\frac{\theta}{2}}\,
\cos\frac{n\theta}{2}.
\]
\end{solution}

\begin{solution}{exo:g12:complex:10}
\emph{1.} $\omega$ est une \omterm{def:g11:quad:discriminant}{racine} $5$-ième de l'unité différente de $1$,
donc la somme des cinq \omterm{def:g11:quad:discriminant}{racines} s'annule
(\cref{thm:g12:complex:rootsofunity}), ce qui s'écrit
$1 + \omega + \omega^2 + \omega^3 + \omega^4 = 0$.

\emph{2.} $u + v = \omega + \omega^2 + \omega^3 + \omega^4 = -1$ par le
point~1. Pour le produit, en utilisant $\omega^5 = 1$~:
\[
uv = (\omega + \omega^4)(\omega^2 + \omega^3)
= \omega^3 + \omega^4 + \omega^6 + \omega^7
= \omega^3 + \omega^4 + \omega + \omega^2 = -1 .
\]
Donc $u, v$ sont les \omterm{def:g11:quad:discriminant}{racines} de $X^2 + X - 1 = 0$~:
$\left\{\frac{-1+\sqrt5}{2}, \frac{-1-\sqrt5}{2}\right\}$. Or
$u = \omega + \conj\omega = 2\cos\frac{2\pi}{5} > 0$ (car
$\frac{2\pi}{5} < \frac{\pi}{2}$), donc $u = \frac{-1 + \sqrt5}{2}$.

\emph{3.} $\cos\frac{2\pi}{5} = \frac u2 = \frac{\sqrt5 - 1}{4}$.
\end{solution}

\begin{solution}{pb:g12:complex:1}
\textbf{1.} $(2 + \iu)(3 - \iu) = 6 - 2\iu + 3\iu + 1 = 7 + \iu$.
\quad $\frac{1 + \iu}{1 - \iu} = \frac{(1 + \iu)^2}{2} =
\frac{2\iu}{2} = \iu$. \quad $\abs{3 + 4\iu} = 5$.

\textbf{2.} $1 + \iu = \sqrt2\,\eu^{\iu\pi/4}$~;
$-2 = 2\eu^{\iu\pi}$~; $1 - \iu\sqrt3 = 2\eu^{-\iu\pi/3}$.

\textbf{3.} $(1 + \iu)^8 = \left(\sqrt2\right)^8
\eu^{8\iu\pi/4} = 16\,\eu^{2\iu\pi} = 16$.

\textbf{4.} $z^2 = \eu^{\iu\pi/2}$~:
$z = \pm\eu^{\iu\pi/4} = \pm\frac{\sqrt2}{2}(1 + \iu)$. Et
$\Delta = 4 - 20 = -16$~: $z = 1 \pm 2\iu$.

\textbf{5.} $z \mapsto \iu z$~: la rotation d'un quart de tour
autour de l'origine. $\abs{z - (1 + \iu)} = 2$~: le cercle de
centre $1 + \iu$ et de rayon $2$.

\textbf{6.} $\left(\omega^k\right)^5 = \eu^{2\iu\pi k} = 1$, et les
cinq points $\eu^{2\iu k\pi/5}$ sont distincts~: ce sont les cinq
solutions de $z^5 = 1$. Elles se placent sur le cercle unité à des
angles égaux de $72^\circ$~: un pentagone régulier ayant un sommet
en $1$.

\textbf{7.} $S = 1 + \omega + \dots + \omega^4$ vérifie
$(1 - \omega)S = 1 - \omega^5 = 0$, et $\omega \neq 1$, donc
$S = 0$. (Le point d'équilibre du pentagone est son centre.)

\textbf{8.} $\omega^4 = \overline\omega$ et
$\omega^3 = \overline{\omega^2}$~: en appariant les \omterm{def:g12:complex:conjugate}{conjugués}, la
\omterm{def:g12:complex:def}{partie réelle} de la somme s'écrit
$1 + 2\cos\frac{2\pi}{5} + 2\cos\frac{4\pi}{5} = 0$.

\textbf{9.} $\cos\frac{4\pi}{5} = 2c^2 - 1$, donc
$1 + 2c + 4c^2 - 2 = 0$, c'est-à-dire $4c^2 + 2c - 1 = 0$~:
$c = \frac{-1 + \sqrt5}{4}$ (la \omterm{def:g11:quad:discriminant}{racine} positive, puisque
$72^\circ$ est aigu)~: $\cos 72^\circ = \frac{\sqrt5 - 1}{4}$.

\textbf{10.} Numériquement $c \approx 0.30902$~: cela correspond à
la calculatrice. Et $\frac{1}{2\varphi} = \frac{1}{1 + \sqrt5} =
\frac{\sqrt5 - 1}{4}$ (après rationalisation)~: donc
$c = \frac{1}{2\varphi}$. Le pentagone est doré de part en part~:
sa diagonale coupe son côté dans le rapport $\varphi$ ---
l'autosimilitude sans fin du pentagramme.

\textbf{11.} Les racines carrées se construisent au compas (la
machine du demi-cercle et de la \omterm{def:g11:stat:mean}{moyenne} géométrique du volume
précédent), si bien qu'un nombre bâti à partir des rationnels par
des racines carrées seulement --- comme $\cos 72^\circ$ --- est
constructible~; les racines cubiques ne le sont pas, et
$\cos 20^\circ$ en réclame une~: le pentagone cède, la trisection
résiste. Le critère de Gauss règle le sort de tout polygone
régulier --- son polygone à 17 côtés, à dix-neuf ans, lui a fait
choisir les mathématiques.

\textbf{12.} $-4 = 4\eu^{\iu\pi}$~: les \omterm{def:g11:quad:discriminant}{racines} quatrièmes ont pour
\omterm{def:g12:complex:modulus}{module} $\sqrt2$ et pour \omterm{def:g12:complex:argument}{arguments}
$\frac{\pi}{4} + k\frac{\pi}{2}$~: ce sont les quatre nombres
$\pm 1 \pm \iu$. En appariant les \omterm{def:g12:complex:conjugate}{conjugués}~:
$(z - 1 - \iu)(z - 1 + \iu) = z^2 - 2z + 2$ et
$(z + 1 - \iu)(z + 1 + \iu) = z^2 + 2z + 2$~: c'est la
factorisation de l'\cref{exo:g12:complex:8}.

\textbf{13.} $(\cos\theta + \iu\sin\theta)^3 = \cos^3\theta +
3\iu\cos^2\theta\sin\theta - 3\cos\theta\sin^2\theta -
\iu\sin^3\theta$~; par la formule de Moivre, cela vaut
$\cos 3\theta + \iu \sin 3\theta$. \omterm{def:g12:complex:def}{Parties réelles}~:
$\cos 3\theta = \cos^3\theta - 3\cos\theta(1 - \cos^2\theta)
= 4\cos^3\theta - 3\cos\theta$ --- et les \omterm{def:g12:complex:def}{parties imaginaires}
livrent $\sin 3\theta$ par-dessus le marché.

\textbf{14.} C'est la \omterm{def:g12:complex:def}{partie réelle} de
$1 + \omega + \dots + \omega^4 = 0$~: la somme vaut $0$ ---
exactement la question 8, vue à travers le noyau de
l'\cref{exo:g12:complex:9}.

\textbf{15.} $8\iu = 8\eu^{\iu\pi/2}$~: \omterm{def:g11:quad:discriminant}{racines}
$2\eu^{\iu\pi/6} = \sqrt3 + \iu$~;
$2\eu^{5\iu\pi/6} = -\sqrt3 + \iu$~;
$2\eu^{3\iu\pi/2} = -2\iu$.

\textbf{16.} $1 + \iu = \sqrt2\,\eu^{\iu\pi/4}$~: l'application
fait tourner le plan de $45^\circ$ et le dilate d'un facteur
$\sqrt2$. Le carré unité devient un carré incliné de côté $\sqrt2$
--- d'aire doublée et pivoté d'un huitième de tour.

\textbf{17.} $\abs{z_1 z_2}^2 = z_1 z_2 \overline{z_1 z_2} =
z_1\overline{z_1}\, z_2\overline{z_2} =
\abs{z_1}^2\abs{z_2}^2$~; il ne reste qu'à prendre les racines
carrées. Puissances de $1 + \iu$~: $1 + \iu$, $2\iu$,
$-2 + 2\iu$, $-4$~: chacune tournée de $45^\circ$ de plus et
allongée d'un facteur $\sqrt2$ --- une spirale sortante de \omterm{def:g12:complex:modulus}{modules}
$\sqrt2, 2, 2\sqrt2, 4, \dots$

\textbf{18.} $\mathrm{j}$ est une racine cubique de l'unité
distincte de $1$, donc (l'\omterm{def:g12:complex:argument}{argument} de la question 7 avec $n = 3$)
$1 + \mathrm{j} + \mathrm{j}^2 = 0$. Critère sur le triangle de
référence~: $a + \mathrm{j}b + \mathrm{j}^2 c = 1 +
\mathrm{j}\cdot\mathrm{j} + \mathrm{j}^2 \cdot \mathrm{j}^2 =
1 + \mathrm{j}^2 + \mathrm{j}^4 = 1 + \mathrm{j}^2 +
\mathrm{j} = 0$~: il est vérifié, comme l'exige le triangle
équilatéral $\left(1, \mathrm{j}, \mathrm{j}^2\right)$.

\textbf{19.} Sur $\C$, un polynôme réel se scinde en facteurs du
premier degré~; ses \omterm{def:g11:quad:discriminant}{racines} non réelles vont par paires \omterm{def:g12:complex:conjugate}{conjuguées}
(conjuguer l'\omterm{def:g10:algebra:equation}{équation}), et chaque paire se multiplie en un trinôme
réel $z^2 - 2\operatorname{Re}(z_0)\,z + \abs{z_0}^2$~: d'où la
factorisation réelle en morceaux de degré $1$ et de degré $2$ ---
la question 12 l'a réalisée sur $z^4 + 4$.

\textbf{20.} Des points~; rotation et dilatation~; des polygones
réguliers~; des identités trigonométriques par développement~; et
un monde clos où toute \omterm{def:g10:algebra:equation}{équation} de degré $n$ possède ses $n$
\omterm{def:g11:quad:discriminant}{racines}. Coda~: à la question 10, le pentagone de $\pi$ a serré la
main du nombre d'or --- les deux hôtes les plus célèbres d'Euclide,
enfin présentés l'un à l'autre par une somme de cinq flèches égale
à zéro.
\end{solution}
